\chapter{Fondements du deep learning}\label{chap:dl}

Ce chapitre rappelle les notions de base nécessaires pour comprendre la suite~: le neurone artificiel, l’apprentissage par descente de gradient, la classification multiclasse et la lutte contre le surapprentissage. Les définitions suivent les ouvrages de référence~\cite{goodfellow2016}.

\section{Intelligence artificielle, apprentissage automatique et deep learning}

L’\textbf{intelligence artificielle} (IA) désigne l’ensemble des techniques qui permettent à une machine de réaliser des tâches demandant habituellement de l’intelligence humaine. L’\textbf{apprentissage automatique} (\emph{machine learning}) en est une branche~: au lieu de programmer des règles à la main, on fournit des exemples et la machine en déduit elle-même une règle. Le \textbf{deep learning} est une branche de l’apprentissage automatique fondée sur des réseaux de neurones comportant de nombreuses couches successives.

Dans notre projet, personne n’écrit de règle du type «~si la peau est jaune avec des taches brunes alors la banane est trop mûre~». Le réseau reçoit des milliers de photographies étiquetées et apprend seul quels indices visuels distinguent les stades de maturité. On parle d’\textbf{apprentissage supervisé}~: chaque image d’entraînement est accompagnée de sa classe correcte (l’\emph{étiquette}). La tâche est une \textbf{classification multiclasse}~: parmi $K=11$ classes possibles, le modèle doit désigner la bonne.

\section{Le neurone artificiel}

Le neurone artificiel est l’unité de calcul de base. Il reçoit plusieurs valeurs d’entrée $x_1,\dots,x_n$, calcule leur somme pondérée augmentée d’un biais, puis applique une fonction d’activation non linéaire $f$ :
\begin{equation}
  y = f\!\left(\sum_{i=1}^{n} w_i x_i + b\right) = f(\mathbf{w}^{\top}\mathbf{x} + b).
\end{equation}
Les \textbf{poids} $w_i$ et le \textbf{biais} $b$ sont les paramètres que le réseau apprend. La figure~\ref{fig:neurone} représente cette opération.

\begin{figure}[H]
\centering
\begin{tikzpicture}[x=1mm,y=1mm]
  \foreach \i/\yy in {1/18,2/0,3/-18}{
    \node[blkw,minimum width=20mm] (x\i) at (0,\yy) {$x_\i$};
    \node[font=\small,navy] at (24,\yy+5) {$w_\i$};
  }
  \node[circle,draw=navy,fill=navy,text=white,minimum size=17mm,font=\large] (s) at (52,0) {$\Sigma + b$};
  \node[blko,minimum width=18mm,minimum height=12mm] (f) at (84,0) {$f(\cdot)$};
  \node[blkg,minimum width=14mm,minimum height=12mm] (y) at (112,0) {$y$};
  \foreach \i in {1,2,3} \draw[arr] (x\i.east) -- (s);
  \draw[arr] (s) -- (f); \draw[arr] (f) -- (y);
  \node[font=\footnotesize,softgray,align=center] at (0,-31) {entrées};
  \node[font=\footnotesize,softgray,align=center] at (52,-17) {somme pondérée\\+ biais};
  \node[font=\footnotesize,softgray,align=center] at (84,-12) {activation};
  \node[font=\footnotesize,softgray,align=center] at (112,-12) {sortie};
\end{tikzpicture}
\caption{Un neurone artificiel~: somme pondérée des entrées, ajout d’un biais, puis fonction d’activation.}\label{fig:neurone}
\end{figure}

Un \textbf{réseau de neurones} organise ces neurones en \textbf{couches}~: chaque couche transforme la sortie de la précédente. Une couche \emph{dense} (ou entièrement connectée) relie chacun de ses neurones à toutes les sorties de la couche précédente. Avec $n$ entrées et $m$ neurones, elle contient $n\,m + m$ paramètres. Un réseau «~profond~» empile de nombreuses couches, chacune construisant des représentations plus abstraites que la précédente.

\section{Fonctions d’activation}

Sans fonction d’activation, empiler des couches revient à composer des applications linéaires, ce qui équivaut à une seule application linéaire~: la profondeur n’apporterait rien. La non-linéarité de $f$ permet au réseau de représenter des fonctions complexes. Le tableau~\ref{tab:activations} présente les fonctions utilisées dans ce rapport.

\begin{table}[H]
\centering
\caption{Principales fonctions d’activation.}\label{tab:activations}
\small
\begin{tabular}{@{}llp{6.2cm}@{}}
\toprule
\textbf{Fonction} & \textbf{Définition} & \textbf{Usage} \\
\midrule
ReLU & $\max(0,x)$ & Couches cachées des CNN~: simple et rapide \\
ReLU6 & $\min(\max(0,x),6)$ & MobileNet~: robuste en basse précision \\
Sigmoïde & $\sigma(x)=\dfrac{1}{1+e^{-x}}$ & Classification binaire \\
Softmax & $\dfrac{e^{z_k}}{\sum_j e^{z_j}}$ & Dernière couche d’un classifieur multiclasse \\
\bottomrule
\end{tabular}
\end{table}

\begin{figure}[H]
\centering
\begin{tikzpicture}
\begin{axis}[width=0.78\textwidth,height=5.4cm,xmin=-4,xmax=8,ymin=-0.5,ymax=7,
  axis lines=middle,xlabel={$x$},ylabel={$f(x)$},xlabel style={at={(ticklabel* cs:1)},anchor=north west},
  ylabel style={at={(ticklabel* cs:1)},anchor=south west},legend style={at={(0.01,0.97)},anchor=north west,font=\footnotesize,draw=none,fill=none},
  grid=both,grid style={gray!15},tick label style={font=\footnotesize},samples=200]
  \addplot[navy,very thick,domain=-4:7]{max(0,x)};\addlegendentry{ReLU}
  \addplot[warnorange,very thick,dashed,domain=-4:8]{min(max(0,x),6)};\addlegendentry{ReLU6}
  \addplot[okgreen,thick,domain=-4:8]{6/(1+exp(-x))};\addlegendentry{Sigmoïde $\times 6$}
\end{axis}
\end{tikzpicture}
\caption{ReLU, ReLU6 et sigmoïde (multipliée par 6 pour la lisibilité).}\label{fig:activations}
\end{figure}

\section{Classification multiclasse~: softmax et entropie croisée}

Pour classer une image en $K$ classes, la dernière couche du réseau produit un vecteur de $K$ scores bruts $z_1,\dots,z_K$ (les \emph{logits}). La fonction \textbf{softmax} les transforme en probabilités~:
\begin{equation}
  p_k = \softmax(\mathbf{z})_k = \frac{e^{z_k}}{\sum_{j=1}^{K} e^{z_j}}, \qquad p_k>0,\quad \sum_{k=1}^{K} p_k = 1.
\end{equation}
Par exemple, avec trois classes et les scores $(4,\,2,\,1)$, on obtient $e^4\approx 54{,}6$, $e^2\approx 7{,}4$, $e^1\approx 2{,}7$, dont la somme vaut $64{,}7$ ; les probabilités sont donc environ $(84\,\%,\ 11\,\%,\ 4\,\%)$. La classe prédite est celle de probabilité maximale, et cette probabilité sert de \textbf{confiance} dans l’application Android (chapitre~\ref{chap:android}).

La qualité des prédictions est mesurée pendant l’entraînement par la fonction de perte \textbf{entropie croisée catégorielle}. Si $y$ est la vraie classe, la perte d’une image vaut
\begin{equation}
  \mathcal{L} = -\log p_{y}.
\end{equation}
Elle est faible quand le modèle attribue une probabilité proche de 1 à la bonne classe, et tend vers l’infini quand il lui en attribue une proche de 0. Quand les classes sont déséquilibrées, on pondère chaque image par un poids $\omega_y$ propre à sa classe~: $\mathcal{L}=-\omega_y\log p_y$ (section~\ref{sec:poids}).

\section{Apprentissage par descente de gradient}

Entraîner un réseau consiste à trouver les paramètres $\theta$ (tous les poids et biais) qui minimisent la perte moyenne sur les données d’entraînement. La méthode standard est la \textbf{descente de gradient}~: à chaque itération, on calcule le gradient $\nabla_\theta\mathcal{L}$ de la perte par rapport aux paramètres, puis on déplace les paramètres dans la direction opposée~:
\begin{equation}
  \theta \leftarrow \theta - \eta\,\nabla_\theta \mathcal{L}.
\end{equation}
Le \textbf{taux d’apprentissage} $\eta$ fixe la taille du pas~: trop grand, l’apprentissage diverge~; trop petit, il stagne. Le gradient est calculé efficacement par l’algorithme de \textbf{rétropropagation}, qui applique la règle de dérivation des fonctions composées de la sortie vers l’entrée du réseau.

En pratique, on ne calcule pas le gradient sur toutes les images à la fois, mais sur de petits \textbf{lots} (\emph{mini-batchs}), de 32 images dans notre projet. Un passage complet sur toutes les images d’entraînement s’appelle une \textbf{époque}. Nous utilisons l’optimiseur \textbf{Adam}~\cite{kingma2015}, qui adapte le pas de chaque paramètre à partir d’estimations des moyennes mobiles du gradient et de son carré~; il est peu sensible au choix de $\eta$ et converge rapidement.

\begin{defbox}{À retenir}
Un réseau apprend en répétant quatre étapes~: (1) \emph{prédire} sur un lot d’images~; (2) \emph{mesurer l’erreur} avec la fonction de perte~; (3) \emph{calculer le gradient} par rétropropagation~; (4) \emph{corriger} les poids avec l’optimiseur.
\end{defbox}

\section{Généralisation et surapprentissage}

Le but n’est pas de bien classer les images d’entraînement, mais de bien classer des images \emph{jamais vues}. Un modèle qui apprend les exemples «~par cœur~» sans dégager de règle générale souffre de \textbf{surapprentissage} (\emph{overfitting})~: son erreur d’entraînement est faible mais son erreur sur de nouvelles données est élevée. Pour le détecter et le limiter, on découpe les données en trois ensembles disjoints~:

\begin{itemize}
  \item l’ensemble d’\textbf{entraînement}, qui sert à ajuster les poids~;
  \item l’ensemble de \textbf{validation}, qui sert à surveiller la généralisation pendant l’entraînement et à choisir les hyperparamètres~;
  \item l’ensemble de \textbf{test}, utilisé une seule fois à la fin pour mesurer les performances finales.
\end{itemize}

Plusieurs techniques de \textbf{régularisation} réduisent le surapprentissage~:

\begin{itemize}
  \item le \textbf{dropout}~\cite{srivastava2014}, qui désactive aléatoirement une fraction des neurones à chaque itération d’entraînement~;
  \item l’\textbf{augmentation de données}, qui génère des variantes des images (rotations, zooms, retournements) pour que le modèle ne dépende pas d’une pose particulière~;
  \item l’\textbf{arrêt anticipé} (\emph{early stopping}), qui interrompt l’entraînement quand la perte de validation cesse de diminuer~;
  \item la \textbf{réduction du taux d’apprentissage sur plateau}, qui diminue $\eta$ quand la validation stagne~;
  \item le \textbf{transfert d’apprentissage}, qui part d’un réseau déjà entraîné (section~\ref{sec:transfer}).
\end{itemize}
